% Phases 0-2 of the xv6-riscv paging research platform.
%
% Build:  pdflatex phases-0-2.tex   (run twice, for the progress bars)
%
% Every number in this deck comes from docs/phase0-phase1-report.md,
% docs/phase2/report.md, docs/phase2/gate-results.txt and
% docs/vm-baseline.txt at tag platform-v1.0.

\documentclass[10pt,aspectratio=169]{beamer}

\usetheme{metropolis}
\usepackage[T1]{fontenc}
\usepackage{lmodern}   % scalable outline fonts; without this, pdflatex
                       % falls back to bitmap EC fonts and every glyph
                       % in the PDF renders blurry
\usepackage{booktabs}
\usepackage{xcolor}

\metroset{sectionpage=progressbar, numbering=fraction, progressbar=frametitle}

% A calmer palette than the default metropolis orange.
\definecolor{accent}{HTML}{2C6E8F}
\definecolor{warn}{HTML}{B5540A}
\definecolor{good}{HTML}{2E7D4F}
\setbeamercolor{alerted text}{fg=warn}
\setbeamercolor{frametitle}{bg=accent}
\setbeamercolor{progress bar}{fg=accent}

\newcommand{\good}[1]{\textcolor{good}{#1}}
\newcommand{\warn}[1]{\textcolor{warn}{#1}}
\newcommand{\ca}{\raisebox{0pt}{$\sim$}}   % "circa", without math spacing

% Every phase frame uses the same four headings, so make them one command.
\newcommand{\phasehead}[2]{%
  \frametitle{#1 \quad\small\normalfont\textcolor{white!75!accent}{#2}}}

\title{Building a Trustworthy Paging Platform in xv6-riscv}
\subtitle{Phases 0--2: tracing, ceilings, and a frozen baseline}
\date{Tag \texttt{platform-v1.0}}

\begin{document}

\maketitle

% ---------------------------------------------------------------- context

\begin{frame}{Where this starts}
  \small
  \textbf{Stock xv6 has no virtual memory management at all.} Pages are
  allocated eagerly, live in RAM until the process exits, and are never
  evicted. No page-fault path for memory, no backing store, no replacement
  policy, no frame metadata.

  \medskip
  \textbf{Before Phase 0, six subsystems had already been built:}
  \begin{itemize}
    \item Lazy allocation (\texttt{sbrklazy}) --- the thing that generates faults
    \item Disk-backed swap outside the filesystem --- raw VirtIO, no journal
    \item A per-frame table --- owner, state, pin count, reference and dirty samples
    \item Three replacement policies --- FIFO, Clock, Aging, switchable at runtime
    \item Prefetch --- synchronous and asynchronous, with the full race matrix
    \item Tracing, and about 35 counters
  \end{itemize}

  \medskip
  \alert{What was missing was not code. It was evidence.} Nothing proved the
  system was correct, and nothing could collect data at the scale a
  machine-learning model needs.
\end{frame}

\begin{frame}{The three phases in one sentence each}
  \small
  \begin{table}
    \begin{tabular}{@{}llp{7.4cm}@{}}
      \toprule
      \textbf{Phase} & \textbf{Name} & \textbf{Purpose} \\
      \midrule
      0 & Trace infrastructure & Make data collection \emph{lossless}, so a model
                                 can never learn from a silently biased sample \\
      \addlinespace
      1 & Raise the ceilings   & Remove the size limits that cap working sets
                                 and block real programs \\
      \addlinespace
      2 & Freeze the platform  & Prove correctness is independent of policy, and
                                 create a fixed point to bisect back to \\
      \bottomrule
    \end{tabular}
  \end{table}

  \medskip
  The order is by dependency, not by calendar. Each phase is useless without
  the one before it: there is no point raising ceilings if the trace is lossy,
  and no point freezing a platform that cannot host a real workload.
\end{frame}

% ================================================================ PHASE 0

\section{Phase 0 --- Trace infrastructure}

\begin{frame}
  \phasehead{Phase 0}{Work done}
  \small
  \begin{itemize}
    \item \textbf{Ring buffer 128 $\to$ 65{,}536 records}; read batch 8 $\to$ 256
    \item \textbf{Record 152\,B $\to$ 64\,B.} Narrowed field widths, and lifted
          the per-record \texttt{version} and \texttt{size} out into a one-time
          header behind a new \texttt{vmtrace\_info()} syscall
    \item \textbf{New \texttt{user/vmdrain.c}} --- a guest-side process that loops
          on \texttt{vmtrace\_read()} and writes raw records to a file
    \item \textbf{New \texttt{tools/decode\_trace.py}} --- host-side decoder and
          validator
    \item \textbf{New \texttt{tools/extract\_file.py}} --- pulls a capture out of
          \texttt{fs.img}, because xv6 cannot hand a file to the host
    \item \textbf{New \texttt{vmctl(VM\_TRACE\_SET\_CAPACITY, n)}} --- shrink the
          ring so overflow tests finish in under a second
    \item \textbf{A drop is fatal, not a warning.} Any capture with
          \texttt{drops != 0} is declared invalid and discarded whole
  \end{itemize}
\end{frame}

\begin{frame}
  \phasehead{Phase 0}{Why}
  \small
  \textbf{The ring held 128 events. Training needs $10^5$ to $10^7$.} Under any
  real workload it wrapped constantly and the drop count climbed without
  bound. Nothing downstream of this works until it is fixed.

  \medskip
  \textbf{Why lossy data is worse than no data:}
  \begin{quote}
    A trace with silent holes lets a model learn from a biased sample with
    nothing flagging it. You then cannot tell ``the model learned something
    real'' from ``the model learned the shape of the ring buffer's overflow
    pattern.''
  \end{quote}

  \medskip
  \textbf{Why shrink the record?} Smaller records mean more events per byte:
  more headroom before the ring overflows, and more events per megabyte of
  capture file.

  \medskip
  \textbf{Why a host-side extractor?} The capture lands inside the guest
  filesystem. Without a way to get it out, the pipeline has no end.
\end{frame}

\begin{frame}
  \phasehead{Phase 0}{Results}
  \small
  \begin{table}
    \begin{tabular}{@{}p{5.4cm}l@{}}
      \toprule
      \textbf{Gate item} & \textbf{Result} \\
      \midrule
      \texttt{vmtest swap-repeat} under tracing
        & \good{27{,}445 records, sequence 1..27445, 0 gaps, 0 drops} \\
      \addlinespace
      Decoder round-trips a 1M-event capture
        & \good{1{,}000{,}000 synthetic records, field-for-field} \\
      \addlinespace
      \dots\ and rejects a broken one
        & \good{catches a removed record \emph{and} a mid-record truncation} \\
      \addlinespace
      Trace buffers are non-pageable
        & \good{4\,MiB in \texttt{.bss}, confirmed with \texttt{nm}} \\
      \bottomrule
    \end{tabular}
  \end{table}

  \medskip
  \textbf{One target was missed, deliberately.} The roadmap asked for
  $\leq 48$\,B per record. Its own list of required fields sums to 73\,B, so 48
  was unreachable without dropping fields it called the minimum set.
  \textbf{64\,B} keeps every field, wastes no padding, and divides the block
  size so a record never straddles a block boundary. Down from 152\,B.
\end{frame}

\begin{frame}
  \phasehead{Phase 0}{Conclusion}
  \small
  \textbf{\good{Claim now licensed:}} \emph{``Every paging event in our
  experiments is recorded --- no sampling, no loss''} --- for any capture that
  passes the check. That one sentence is what makes every later number
  defensible.

  \medskip
  \textbf{Ruled out:} that a model was trained on a silently biased subset of
  events.

  \medskip
  \textbf{What the drop mechanism actually does} (learned by getting a test
  wrong): on overrun the kernel evicts the oldest record \emph{and} rewrites
  the incoming one as a \texttt{DROP} marker. Sustained overrun therefore
  turns the entire surviving window into markers.
  \begin{itemize}
    \item \good{Good:} an overrun window can never be mistaken for data,
          whichever end you read it from.
    \item \warn{Cost:} an overrun capture is not partially salvageable --- you
          lose the window, not just the oldest record.
  \end{itemize}
  The right trade for training data, but a conscious one.

  \medskip
  \textbf{Still cannot say:} anything about replacement quality. This phase
  concerns instrument integrity only.
\end{frame}

% ================================================================ PHASE 1

\section{Phase 1 --- Raise the ceilings}

\begin{frame}
  \phasehead{Phase 1}{Work done}
  \small
  \begin{table}
    \begin{tabular}{@{}llll@{}}
      \toprule
      \textbf{Constant} & \textbf{Was} & \textbf{Now} & \textbf{Effect} \\
      \midrule
      \texttt{NSWAPSLOTS} & 1024    & 8192      & swap 4\,MB $\to$ \good{32\,MiB} \\
      \texttt{USERSTACK}  & 1       & 16        & user stack 4\,KB $\to$ \good{64\,KiB} \\
      \texttt{FSSIZE}     & 2000    & 100000    & filesystem 2\,MB $\to$ \good{100\,MB} \\
      \texttt{MAXFILE}    & 268 blk & 65803 blk & max file 274\,KB $\to$ \good{64.26\,MiB} \\
      \texttt{NBUF}       & 30      & 30        & \warn{deliberately unchanged} \\
      \bottomrule
    \end{tabular}
  \end{table}

  \vspace{-2mm}
  \begin{itemize}
    \item \texttt{MAXFILE} grew by adding a \textbf{doubly-indirect block level}
          (\texttt{NDIRECT} 12 $\to$ 11 plus one doubly-indirect entry), and
          extending \texttt{bmap()} and \texttt{itrunc()} to match
    \item \textbf{The same change had to be made again in \texttt{mkfs/mkfs.c}},
          which keeps its own independent copy of the block-mapping logic
    \item New \texttt{user/bigfiletest.c} and \texttt{tools/fsck\_xv6.py},
          a host-side leak and consistency checker
    \item \texttt{mkfs} now truncates the image instead of writing 100{,}000
          zero blocks --- rebuild dropped from \ca 90\,s to seconds
  \end{itemize}
\end{frame}

\begin{frame}
  \phasehead{Phase 1}{Why}
  \small
  \textbf{Every limit here blocks something specific downstream.}

  \medskip
  \begin{itemize}
    \item \textbf{4\,MB swap} caps working sets at 4\,MB. A reviewer can
          reasonably ask whether anything observed at that size generalises
          past cache size.
    \item \textbf{274\,KB maximum file} is smaller than any real program binary
          or input dataset. It rules out hosting third-party software
          (Phase 3) and workloads driven by real input data (Phase 4)
          outright.
    \item \textbf{4\,KB user stack} means ordinary recursion faults.
    \item \textbf{\texttt{NBUF} left alone on purpose.} Growing the block cache
          would move the interesting behaviour out of the paging layer ---
          which is exactly where the research needs it to stay.
  \end{itemize}

  \medskip
  \textbf{The risk that was named in advance:} the \texttt{mkfs} half is the
  one that gets forgotten. The symptom would be a kernel that can \emph{read}
  large files and a build that cannot \emph{create} them.
\end{frame}

\begin{frame}
  \phasehead{Phase 1}{Results}
  \small
  \begin{itemize}
    \item \good{\texttt{usertests -q} passes. \texttt{grind} passes.}
    \item \good{A 5\,MiB file} writes, reads back byte-exact, and deletes with
          no block leak --- confirmed by an independent host-side fsck
          (blocks used $=$ blocks reachable, 0 leaked)
    \item \good{\texttt{vmtest all} and \texttt{prefetchtest all} pass under
          FIFO, Clock and Aging}, in both the release and the \texttt{VM\_DEBUG}
          build
    \item \good{\texttt{vmcheck} clean after every one}
  \end{itemize}

  \medskip
  \alert{Ten distinct failures on the way --- and five of them shared a single
  root cause.}

  \medskip
  \texttt{USERSTACK} 1 $\to$ 16 added fifteen pages to \emph{every} process's
  baseline resident set: from about 11 frames to 27. Five separate tests had
  encoded ``the thing I am testing is much larger than everything else this
  process holds'' as a \textbf{fixed page count}. One constant in
  \texttt{param.h} inverted that relationship across two test programs and
  five tests --- \textbf{none of which mentioned \texttt{USERSTACK}}.
\end{frame}

\begin{frame}
  \phasehead{Phase 1}{Conclusion}
  \small
  \textbf{\good{Claim now licensed:}} \emph{``Working-set sizes in this study
  are bounded by experimental design, not by an artefact of the teaching
  kernel.''}

  \medskip
  \textbf{Ruled out:} \emph{``your policy differences are just noise, because
  nothing ever exceeded a few megabytes.''}
  \quad\textbf{Unlocked:} the ability to host a real third-party binary at all
  (Phase 3), and workloads whose data does not fit in the block cache
  (Phase 4).

  \medskip
  \textbf{The transferable lesson, from the five-in-one failure:}
  \begin{quote}
    Derive test magnitudes from \emph{observed state}, never from literals.
    And where a test's meaning depends on a state it establishes indirectly
    --- through a limit, a policy, a fault --- \textbf{assert that the state was
    actually established.}
  \end{quote}

  \medskip
  \warn{A caution that outlives this phase:} every process now carries about
  27 cold resident pages before it allocates anything. A naive
  ``limit $=$ N'' experiment is therefore not measuring ``a working set of N
  pages'' --- it measures N minus a policy-dependent slice of process overhead.
\end{frame}

% ================================================================ PHASE 2

\section{Phase 2 --- Freeze the platform}

\begin{frame}
  \phasehead{Phase 2}{Work done}
  \small
  \begin{itemize}
    \item \textbf{Audited every PTE mutation in the kernel} --- 24 sites, each
          classified as page-table construction, mapping installation, A/D
          sampling, swap transition, permission change, or cleanup
    \item \textbf{New \texttt{vmtest data-invariance}} --- one deterministic
          workload run under all 3 policies $\times$ 2 prefetch modes,
          requiring byte-exact agreement between them
    \item \textbf{New \texttt{vmtest baseline}} --- the same workload with
          prefetch off, printing the reference numbers to bisect against
    \item \textbf{Provenance on every transcript} --- a boot banner printing the
          compile-time paging configuration, plus a harness header carrying the
          commit, the toolchain versions and the build configuration
    \item \textbf{New \texttt{tools/phase2\_gate.sh}} --- the matrix driver
    \item \textbf{Ran the whole matrix:} debug and release builds, all three
          policies, ten randomised seeds each, \texttt{CPUS=3} stress, and the
          trace-integrity checks
  \end{itemize}
\end{frame}

\begin{frame}
  \phasehead{Phase 2}{Why}
  \small
  \textbf{Everything after this point is research.} When a later experiment
  produces an implausible number, this tag is the point you bisect back to.
  Skipping it means later debugging has no fixed point.

  \medskip
  \textbf{The specific disaster being insured against:}
  \begin{quote}
    Discovering at the end of the thesis that a reported improvement came from
    a lifecycle bug --- a leaked swap slot, a lost dirty bit, an eviction that
    quietly dropped data --- rather than from better victim selection.
  \end{quote}

  \medskip
  \textbf{Why a new test was needed.} The requirement ``all policies produce
  identical application data'' had \emph{no direct test}. Every existing policy
  test fixes one policy and checks that process alone --- which establishes
  that each policy is individually correct, but never that they \emph{agree
  with each other}.

  \medskip
  \textbf{Why print the configuration.} A number is only evidence if you can
  say which kernel, which commit and which settings produced it --- recorded at
  the time, not reconstructed from memory afterwards.
\end{frame}

\begin{frame}
  \phasehead{Phase 2}{Results --- the matrix}
  \begin{columns}[T]
    \begin{column}{0.44\textwidth}
      \small
      \textbf{\good{74 steps. 0 failures.}}
      \medskip
      \begin{itemize}\small
        \item \texttt{CPUS=1} release
        \item \texttt{CPUS=1 VM\_DEBUG=1}
        \item \texttt{CPUS=3}, stress only
        \item FIFO, Clock, Aging --- each under both prefetch modes
        \item 10 randomised seeds per configuration
        \item \texttt{usertests -q}, \texttt{grind}, big-file and fsck checks
        \item about 80 minutes end to end
      \end{itemize}
      \medskip
      All 15 acceptance criteria ticked with named evidence.
    \end{column}
    \begin{column}{0.54\textwidth}
      \small\textbf{The data-invariance result}\\[4pt]
      \scriptsize
      \begin{tabular}{@{}llrrrl@{}}
        \toprule
        Policy & Mode & Evict & Faults & Waste & Checksum \\
        \midrule
        FIFO  & sync  & 1567 & 1039 &  48 & \good{AB0C8578} \\
        FIFO  & async & 1567 & 1039 &  48 & \good{AB0C8578} \\
        Clock & sync  & 1746 & 1127 & 347 & \good{AB0C8578} \\
        Clock & async & 1721 & 1102 & 325 & \good{AB0C8578} \\
        Aging & sync  & 1489 & 1000 &  14 & \good{AB0C8578} \\
        Aging & async & 1487 &  998 &  12 & \good{AB0C8578} \\
        \bottomrule
      \end{tabular}

      \medskip
      \small
      17\% spread in evictions, 13\% in faults, a factor of \textbf{29} in
      wasted prefetches --- and \textbf{identical data every time}.
    \end{column}
  \end{columns}
\end{frame}

\begin{frame}
  \phasehead{Phase 2}{Results --- three defects found}
  \small
  \textbf{1. A real kernel bug: async prefetch reads billed to the wrong
  process.}\\
  \texttt{account\_io()} took its process from \texttt{myproc()}, but the async
  prefetch worker is a \emph{separate kernel process}. Its reads went to the
  worker's own counters, which nothing reads. An owner's \texttt{page\_reads}
  came out \emph{exactly} equal to its \texttt{swap\_faults} --- an
  \alert{41\% undercount of device traffic}. Phase 5 compares I/O across
  prefetch modes; it would have concluded async does far less I/O than sync.
  The truth is the opposite.

  \medskip
  \textbf{2. A test that had silently stopped testing.}\\
  \texttt{vmtest random} used a fixed 24-page region against a limit of 35. It
  paged eighteen times during warm-up, then ran \textbf{400{,}000 operations
  without a single fault} --- and passed, every time, for the life of the
  project. Now 51 pages against 35: \good{37{,}495 swap faults per 100k
  operations}.

  \medskip
  \textbf{3. An artefact in Phase 2's own new code.}\\
  Reconfiguring one process in a loop meant FIFO ran at limit 34 and Clock at
  12. It produced a \alert{publishable-looking ``FIFO evicts 34\% fewer
  pages.''} Fixed with a forked child per configuration and the limit computed
  once. Corrected: Clock is \good{8.0\%} better than FIFO.
\end{frame}

\begin{frame}
  \phasehead{Phase 2}{Results --- one measured negative result}
  \small
  \textbf{Lossless tracing does not yet reach dataset scale.} Fixing defect 2
  raised the fault rate to a realistic level, and the ring could not keep up.

  \smallskip
  \begin{table}\small
    \begin{tabular}{@{}lr@{}}
      \toprule
      Events emitted (deterministic across runs)   & 302{,}243 \\
      Records that reached the file                & 161{,}285 \quad 53.4\% \\
      \dots\ of which are \texttt{DROP} markers, not data & 65{,}588 \\
      \midrule
      \textbf{Arrived with an intact event type}   & \alert{95{,}697 \quad 31.7\%} \\
      \midrule
      Emission / drain rate during the run & \ca 1{,}980 / \ca 1{,}050 per s \\
      Wall-clock cost of tracing           & 72.5\,s $\to$ 152.4\,s (2.1x) \\
      \bottomrule
    \end{tabular}
  \end{table}

  \vspace{-1mm}
  \textbf{This is a finding, not a gap.} Phase 0's gate explicitly allowed
  for it. The matrix now records \emph{both} outcomes on purpose --- one
  lossless capture, and one that \good{must be rejected} --- which is
  positive evidence that the drop discipline fires.

  \smallskip
  \textbf{Three remedies, costed:} pace the workload (changes nothing about
  the system under test); enlarge the ring (bigger backlog, but shifts every
  absolute page count); or emit fewer event types (loses I/O latency).
\end{frame}

\begin{frame}
  \phasehead{Phase 2}{Conclusion}
  \small
  \textbf{\good{Claim now licensed --- and it is the load-bearing one:}}
  \begin{quote}
    \emph{The platform is correct, and correctness is independent of
    replacement policy.}
  \end{quote}
  Six configurations differ by 17\% in evictions and by a factor of 29 in
  wasted prefetches, yet all six produce the same byte-exact result --- equal to
  a checksum computed from the pattern definition \emph{without touching memory
  at all}. \textbf{So any later difference between policies is a performance
  difference, not a bug.}

  \medskip
  \textbf{Reproducibility, measured rather than assumed:} synchronous runs are
  bit-identical across builds and across runs; asynchronous runs reproduce in
  \emph{data} only; and the counters belong to a \emph{build of the test
  binary} --- editing the test file moved Clock's baseline by two faults with no
  change to policy or workload.

  \medskip
  \warn{What this does \emph{not} license:} nothing about replacement quality
  (one workload, one capacity), and nothing about multicore --- no cross-hart
  TLB shootdown exists, so \texttt{CPUS=3} is stress only.
\end{frame}

% ================================================================ COMBINED

\section{Taken together}

\begin{frame}{What the three phases jointly establish}
  \small
  \begin{enumerate}
    \item \textbf{The instrument is honest} \emph{(Phase 0)}.\\
      Every recorded event is recorded, and a capture that lost events says so
      loudly and is thrown away whole rather than quietly truncated.

    \smallskip
    \item \textbf{The instrument has range} \emph{(Phase 1)}.\\
      32\,MiB of swap, 64\,MiB files, 64\,KiB stacks. Results cannot be
      dismissed as artefacts of a 4\,MB ceiling, and real programs with real
      input data now fit.

    \smallskip
    \item \textbf{The instrument is calibrated} \emph{(Phase 2)}.\\
      Correctness is provably independent of policy, so later differences
      between policies are performance --- and there is a tagged fixed point
      with recorded baselines to bisect against.
  \end{enumerate}

  \medskip
  \begin{quote}
    The project moved from ``a paging system exists'' to ``a paging system
    whose measurements can be trusted.''
  \end{quote}

  \textbf{None of that is a research result. All of it is the precondition for
  one.}
\end{frame}

\begin{frame}{The one lesson that kept recurring}
  \small
  \begin{quote}
    \large
    \textbf{A fixed magnitude competing against a baseline that grew
    underneath it.}
  \end{quote}

  \medskip
  \texttt{USERSTACK} 1 $\to$ 16 raised every process's post-\texttt{exec}
  resident set from about 11 frames to 27. That single change broke
  \textbf{five tests in Phase 0/1} and \textbf{a sixth in Phase 2} --- and
  \alert{every one of them kept passing} while testing less and less.

  \medskip
  \textbf{The rule that follows, and it applies to every workload built from
  here on:}
  \begin{itemize}
    \item Derive region sizes and resident limits from \emph{observed} state,
          never from literals.
    \item Assert afterwards that the setup actually took effect --- that the
          workload really did fault and evict.
  \end{itemize}

  \medskip
  This matters most in Phase 4: a benchmark whose working set accidentally fits
  the resident limit will produce a beautiful, meaningless reuse-distance
  histogram.
\end{frame}

\begin{frame}{The honest ledger}
  \small
  \begin{table}\small
    \begin{tabular}{@{}lp{8.4cm}@{}}
      \toprule
      \textbf{13} & defects found across the three phases
                    (10 in Phase 0/1, 3 in Phase 2) \\
      \addlinespace
      \textbf{2}  & of the Phase 2 defects would have produced plausible,
                    \emph{publishable-looking} numbers pointing the wrong way \\
      \addlinespace
      \textbf{0}  & of those two were visible from a test result --- both were
                    found by looking at a column that was not the answer \\
      \addlinespace
      \textbf{1}  & blocker carried forward, with a measured number and three
                    costed remedies \\
      \bottomrule
    \end{tabular}
  \end{table}

  \vspace{-1mm}
  \begin{quote}
    \textbf{The freeze did not confirm that the platform was correct. It found
    the three places where it was not.}
  \end{quote}

  \medskip
  That is the argument for these phases existing, and it is worth making in
  those terms: their value is measured in what they caught, not in the green
  checkmarks they produced.
\end{frame}

\begin{frame}{What comes next}
  \small
  \textbf{Immediately:}
  \begin{itemize}
    \item \textbf{Phase 3} --- a hard-timeboxed attempt to run SQLite. It buys a
          \emph{methodological independence} argument (``software we did not
          write''), not a capability. Abandonable by design; failing is itself
          a reportable systems finding.
    \item \textbf{Phase 4} --- six workloads with deliberately distinct
          reuse-distance profiles. \alert{Must resolve the tracing rate limit
          first.}
  \end{itemize}

  \medskip
  \textbf{Then the first genuine research result:}
  \begin{itemize}
    \item \textbf{Phase 5} --- a Belady optimal oracle computed offline, and the
          measured gap between Clock and optimal. \textbf{This result exists
          whether or not the machine learning works.} If the gap is small
          everywhere, the honest conclusion is ``learned replacement is not
          worth its cost in this regime'' --- which is a real contribution, not
          a failure.
  \end{itemize}

  \vspace{-1mm}
  \centering
  \textcolor{gray}{Every phase from 3 onward has a defensible negative outcome.
  That property is deliberate, and it is the reason for this ordering.}
\end{frame}

\begin{frame}[standout]
  Phases 0--2 complete\\[6pt]
  {\large\normalfont Tag \texttt{platform-v1.0} --- 74 steps, 0 failures}
\end{frame}

\end{document}
